\documentclass[11pt,a4paper]{article}

\usepackage[T1]{fontenc}
\usepackage[utf8]{inputenc}
\usepackage[brazil]{babel}
\usepackage[a4paper, margin=2.2cm]{geometry}
\usepackage{xcolor}
\usepackage{amssymb}
\usepackage{listings}
\usepackage{enumitem}
\usepackage{booktabs}
\usepackage{titlesec}
\usepackage{tcolorbox}
\usepackage{fancyhdr}
\usepackage[colorlinks=true, linkcolor=azul, urlcolor=azul]{hyperref}

\definecolor{azul}{HTML}{1F4E8B}
\definecolor{cinzaCode}{HTML}{F4F4F4}
\definecolor{verdeOk}{HTML}{1B7F3A}
\definecolor{laranjaWarn}{HTML}{B86A00}
\definecolor{vermelhoErr}{HTML}{B23A48}
\definecolor{comentario}{HTML}{6A737D}
\definecolor{keyword}{HTML}{0033B3}
\definecolor{string}{HTML}{B23A48}

\lstdefinelanguage{JavaScript}{
  keywords={break, case, catch, continue, debugger, default, delete, do, else, false, finally, for, function, if, in, instanceof, new, null, return, switch, this, throw, true, try, typeof, var, void, while, with, let, const, async, await, of, import, export, from, class, extends, super},
  morekeywords=[2]{document, window, console},
  sensitive=true,
  morecomment=[l]{//},
  morecomment=[s]{/*}{*/},
  morestring=[b]',
  morestring=[b]",
  morestring=[b]`,
}

\lstdefinelanguage{HTML5}{
  language=HTML,
  morekeywords={meta,charset,article,header,main,section,nav,aside,footer,figure,figcaption},
  alsoletter={-},
}

\lstdefinelanguage{CSS}{
  morekeywords={color,background,font,padding,margin,border,width,height,display,position,top,left,right,bottom,flex,grid,box-sizing,border-radius,box-shadow,text-align,font-family,font-size,font-weight,cursor},
  sensitive=false,
  morecomment=[s]{/*}{*/},
  morestring=[b]',
  morestring=[b]",
  alsoletter={-},
}

\lstdefinestyle{codigo}{
  backgroundcolor=\color{cinzaCode},
  basicstyle=\ttfamily\footnotesize,
  keywordstyle=\color{keyword}\bfseries,
  commentstyle=\color{comentario}\itshape,
  stringstyle=\color{string},
  numbers=left,
  numberstyle=\tiny\color{gray},
  numbersep=8pt,
  frame=single,
  rulecolor=\color{gray!40},
  framesep=6pt,
  breaklines=true,
  showstringspaces=false,
  tabsize=2,
  captionpos=b,
  literate=
    {á}{{\'a}}1 {à}{{\`a}}1 {ã}{{\~a}}1 {â}{{\^a}}1 {ä}{{\"a}}1
    {é}{{\'e}}1 {ê}{{\^e}}1 {í}{{\'i}}1 {ó}{{\'o}}1 {ô}{{\^o}}1
    {õ}{{\~o}}1 {ú}{{\'u}}1 {ç}{{\c{c}}}1
    {Á}{{\'A}}1 {É}{{\'E}}1 {Í}{{\'I}}1 {Ó}{{\'O}}1 {Ú}{{\'U}}1
    {Ç}{{\c{C}}}1
}
\lstset{style=codigo}

\newtcolorbox{caixa-dica}{
  colback=blue!4, colframe=azul, boxrule=0.5pt, arc=2pt,
  left=8pt, right=8pt, top=4pt, bottom=4pt,
  title=\textbf{Dica}, fonttitle=\color{white}
}

\newtcolorbox{caixa-atencao}{
  colback=orange!6, colframe=laranjaWarn, boxrule=0.5pt, arc=2pt,
  left=8pt, right=8pt, top=4pt, bottom=4pt,
  title=\textbf{Atenção}, fonttitle=\color{white}
}

\newtcolorbox{caixa-pratica}{
  colback=green!4, colframe=verdeOk, boxrule=0.5pt, arc=2pt,
  left=8pt, right=8pt, top=4pt, bottom=4pt,
  title=\textbf{Pratique}, fonttitle=\color{white}
}

\titleformat{\section}{\Large\bfseries\color{azul}}{\thesection}{0.6em}{}
\titleformat{\subsection}{\large\bfseries\color{azul!85!black}}{\thesubsection}{0.6em}{}

\pagestyle{fancy}
\fancyhf{}
\fancyhead[L]{\small\color{azul}DSW --- CRUD de Produtos}
\fancyhead[R]{\small UNEMAT 2026.1}
\fancyfoot[C]{\small\thepage}
\renewcommand{\headrulewidth}{0.4pt}

\title{\vspace{-1.5cm}
  \color{azul}\Huge Tutorial Prático\\
  \LARGE CRUD de Produtos com Express, Bun e MySQL
}
\author{\large Prof. Ivan Luiz Pedroso Pires\\\small Desenvolvimento de Software para Web --- UNEMAT}
\date{\large 01 de junho de 2026}

\begin{document}

\maketitle
\thispagestyle{fancy}

\begin{abstract}
\noindent Neste tutorial você vai construir uma aplicação \textbf{CRUD} (Create, Read, Update, Delete) completa de produtos. O \textbf{backend} usa \texttt{Express.js} rodando no \texttt{Bun}, com banco \texttt{MySQL} acessado pelo \texttt{mysql2}, organizado no padrão \textbf{Routes e Controllers}. O \textbf{frontend} é \texttt{HTML} puro com \texttt{JavaScript} usando \texttt{fetch}. Duração estimada: \textbf{2h30}.
\end{abstract}

\tableofcontents
\vspace{0.5cm}
\hrule
\vspace{0.5cm}

% ============================================================
\section{Pré-requisitos}
% ============================================================

Antes de começar, confirme que você tem instalado:

\begin{itemize}[leftmargin=*, itemsep=2pt]
  \item \textbf{Bun} (\url{https://bun.sh}) --- runtime JavaScript
  \item \textbf{MySQL Server} rodando em \texttt{localhost:3306}
  \item \textbf{Editor de código} (VS Code recomendado)
  \item \textbf{Navegador} moderno (Chrome, Firefox ou Edge)
  \item Cliente HTTP opcional: \textbf{Insomnia} ou \textbf{Thunder Client}
\end{itemize}

Teste no terminal:

\begin{lstlisting}[language=bash]
bun --version
mysql --version
\end{lstlisting}

\begin{caixa-atencao}
Se algum comando não for reconhecido, instale antes de continuar. Bun: \texttt{curl -fsSL https://bun.sh/install | bash}.
\end{caixa-atencao}

% ============================================================
\section{Visão geral do projeto}
% ============================================================

Vamos construir duas partes que se comunicam por HTTP:

\begin{center}
\begin{tabular}{ll}
\toprule
\textbf{Frontend (porta livre)} & \textbf{Backend (porta 3000)} \\
\midrule
\texttt{index.html}            & \texttt{server.js} \\
\texttt{style.css}             & \texttt{routes/produtoRoutes.js} \\
\texttt{script.js} (fetch)     & \texttt{controllers/produtoController.js} \\
                               & \texttt{config/database.js} \\
\bottomrule
\end{tabular}
\end{center}

\textbf{Fluxo de uma requisição}:

\begin{center}
\texttt{Browser} $\rightarrow$ \texttt{fetch()} $\rightarrow$ \texttt{Express} $\rightarrow$ \texttt{Route} $\rightarrow$ \texttt{Controller} $\rightarrow$ \texttt{MySQL}
\end{center}

\textbf{Estrutura final de pastas}:

\begin{lstlisting}[language=bash]
01-06-2025/
|-- backend/
|   |-- database.sql
|   |-- package.json
|   `-- src/
|       |-- config/database.js
|       |-- controllers/produtoController.js
|       |-- routes/produtoRoutes.js
|       `-- server.js
`-- frontend/
    |-- index.html
    |-- style.css
    `-- script.js
\end{lstlisting}

% ============================================================
\section{Passo 1 --- Criar o banco de dados}
% ============================================================

Acesse o MySQL e crie a tabela \texttt{produtos}.

\begin{lstlisting}[language=SQL, caption={\texttt{backend/database.sql}}]
CREATE DATABASE IF NOT EXISTS dsw_crud;
USE dsw_crud;

CREATE TABLE IF NOT EXISTS produtos (
  id INT AUTO_INCREMENT PRIMARY KEY,
  nome VARCHAR(100) NOT NULL,
  preco DECIMAL(10,2) NOT NULL,
  estoque INT NOT NULL DEFAULT 0,
  criado_em TIMESTAMP DEFAULT CURRENT_TIMESTAMP
);

INSERT INTO produtos (nome, preco, estoque) VALUES
  ('Caneta Azul', 2.50, 100),
  ('Caderno 96 folhas', 15.90, 50),
  ('Borracha', 1.20, 200);
\end{lstlisting}

Execute o script:

\begin{lstlisting}[language=bash]
mysql -u root -p < backend/database.sql
\end{lstlisting}

Valide criando a tabela:

\begin{lstlisting}[language=SQL]
USE dsw_crud;
SELECT * FROM produtos;
\end{lstlisting}

\begin{caixa-dica}
A coluna \texttt{id} é \texttt{AUTO\_INCREMENT}: o MySQL gera o número sozinho. Nunca passe \texttt{id} no \texttt{INSERT}.
\end{caixa-dica}

% ============================================================
\section{Passo 2 --- Inicializar o backend com Bun}
% ============================================================

Crie as pastas e o \texttt{package.json}:

\begin{lstlisting}[language=bash]
mkdir -p backend/src/{config,controllers,routes}
cd backend
bun init -y
\end{lstlisting}

Edite o \texttt{package.json} ficando assim:

\begin{lstlisting}[language=JavaScript, caption={\texttt{backend/package.json}}]
{
  "name": "backend",
  "version": "1.0.0",
  "type": "module",
  "scripts": {
    "dev": "bun --watch src/server.js",
    "start": "bun src/server.js"
  },
  "dependencies": {
    "cors": "^2.8.5",
    "express": "^4.21.2",
    "mysql2": "^3.11.5"
  }
}
\end{lstlisting}

Instale as dependências:

\begin{lstlisting}[language=bash]
bun install
\end{lstlisting}

\begin{caixa-atencao}
O campo \texttt{"type": "module"} permite usar \texttt{import/export} em vez de \texttt{require}.
\end{caixa-atencao}

% ============================================================
\section{Passo 3 --- Conexão com o MySQL}
% ============================================================

Crie um \textbf{pool de conexões} (mais eficiente que abrir/fechar a cada query).

\begin{lstlisting}[language=JavaScript, caption={\texttt{backend/src/config/database.js}}]
import mysql from 'mysql2/promise';

const pool = mysql.createPool({
  host: 'localhost',
  port: 3306,
  user: 'root',
  password: '',
  database: 'dsw_crud',
  waitForConnections: true,
  connectionLimit: 10,
});

export default pool;
\end{lstlisting}

\begin{caixa-dica}
Use \texttt{mysql2/promise}: permite \texttt{await} nas queries em vez de callbacks aninhados.
\end{caixa-dica}

% ============================================================
\section{Passo 4 --- Controller de Produtos}
% ============================================================

O \textbf{controller} contém a lógica do CRUD. Cada função recebe \texttt{(req, res)} e responde JSON.

\begin{lstlisting}[language=JavaScript, caption={\texttt{backend/src/controllers/produtoController.js (parte 1)}}]
import db from '../config/database.js';

export const listar = async (req, res) => {
  try {
    const [rows] = await db.query(
      'SELECT * FROM produtos ORDER BY id DESC'
    );
    res.json(rows);
  } catch (error) {
    res.status(500).json({ erro: error.message });
  }
};

export const buscarPorId = async (req, res) => {
  try {
    const { id } = req.params;
    const [rows] = await db.query(
      'SELECT * FROM produtos WHERE id = ?', [id]
    );
    if (rows.length === 0) {
      return res.status(404).json({ erro: 'Produto nao encontrado' });
    }
    res.json(rows[0]);
  } catch (error) {
    res.status(500).json({ erro: error.message });
  }
};
\end{lstlisting}

\begin{lstlisting}[language=JavaScript, caption={\texttt{produtoController.js (parte 2 --- escrita)}}]
export const criar = async (req, res) => {
  try {
    const { nome, preco, estoque } = req.body;
    if (!nome || preco == null) {
      return res.status(400).json({ erro: 'Nome e preco obrigatorios' });
    }
    const [result] = await db.query(
      'INSERT INTO produtos (nome, preco, estoque) VALUES (?, ?, ?)',
      [nome, preco, estoque ?? 0]
    );
    res.status(201).json({
      id: result.insertId, nome, preco, estoque: estoque ?? 0
    });
  } catch (error) {
    res.status(500).json({ erro: error.message });
  }
};

export const atualizar = async (req, res) => {
  try {
    const { id } = req.params;
    const { nome, preco, estoque } = req.body;
    const [result] = await db.query(
      'UPDATE produtos SET nome=?, preco=?, estoque=? WHERE id=?',
      [nome, preco, estoque, id]
    );
    if (result.affectedRows === 0) {
      return res.status(404).json({ erro: 'Produto nao encontrado' });
    }
    res.json({ id: Number(id), nome, preco, estoque });
  } catch (error) {
    res.status(500).json({ erro: error.message });
  }
};

export const deletar = async (req, res) => {
  try {
    const { id } = req.params;
    const [result] = await db.query(
      'DELETE FROM produtos WHERE id = ?', [id]
    );
    if (result.affectedRows === 0) {
      return res.status(404).json({ erro: 'Produto nao encontrado' });
    }
    res.json({ mensagem: 'Produto removido com sucesso' });
  } catch (error) {
    res.status(500).json({ erro: error.message });
  }
};
\end{lstlisting}

\begin{caixa-atencao}
\textbf{Sempre use \texttt{?} (prepared statements)} para inserir variáveis nas queries. Concatenação direta abre brecha para \textbf{SQL Injection}.
\end{caixa-atencao}

% ============================================================
\section{Passo 5 --- Rotas de Produtos}
% ============================================================

As \textbf{rotas} só declaram qual URL e método HTTP chamam qual função do controller.

\begin{lstlisting}[language=JavaScript, caption={\texttt{backend/src/routes/produtoRoutes.js}}]
import { Router } from 'express';
import {
  listar,
  buscarPorId,
  criar,
  atualizar,
  deletar,
} from '../controllers/produtoController.js';

const router = Router();

router.get('/', listar);
router.get('/:id', buscarPorId);
router.post('/', criar);
router.put('/:id', atualizar);
router.delete('/:id', deletar);

export default router;
\end{lstlisting}

\textbf{Mapa de endpoints} (combinando o prefixo definido no \texttt{server.js}):

\begin{center}
\begin{tabular}{lll}
\toprule
\textbf{Método} & \textbf{URL} & \textbf{Função} \\
\midrule
GET    & \texttt{/produtos}     & listar todos \\
GET    & \texttt{/produtos/:id} & buscar um \\
POST   & \texttt{/produtos}     & criar \\
PUT    & \texttt{/produtos/:id} & atualizar \\
DELETE & \texttt{/produtos/:id} & deletar \\
\bottomrule
\end{tabular}
\end{center}

% ============================================================
\section{Passo 6 --- Servidor Express}
% ============================================================

O \texttt{server.js} amarra tudo: aplica middlewares, registra rotas, sobe o servidor.

\begin{lstlisting}[language=JavaScript, caption={\texttt{backend/src/server.js}}]
import express from 'express';
import cors from 'cors';
import produtoRoutes from './routes/produtoRoutes.js';

const app = express();

app.use(cors());
app.use(express.json());

app.use('/produtos', produtoRoutes);

app.get('/', (req, res) => {
  res.json({ mensagem: 'API do CRUD de Produtos rodando' });
});

const PORT = 3000;
app.listen(PORT, () => {
  console.log(`Servidor rodando em http://localhost:${PORT}`);
});
\end{lstlisting}

Rode o servidor:

\begin{lstlisting}[language=bash]
bun run dev
\end{lstlisting}

Você deve ver: \texttt{\color{verdeOk}Servidor rodando em http://localhost:3000}.

\begin{caixa-dica}
\textbf{Middlewares importam}:
\begin{itemize}
  \item \texttt{cors()}: libera chamadas do frontend (origem diferente).
  \item \texttt{express.json()}: faz o parse do corpo JSON das requisições.
\end{itemize}
A \textbf{ordem} é importante: middlewares antes das rotas.
\end{caixa-dica}

% ============================================================
\section{Passo 7 --- Testar a API}
% ============================================================

Antes de codar o frontend, valide a API. No Insomnia/Thunder Client:

\begin{enumerate}[itemsep=2pt]
  \item \textbf{GET} \texttt{http://localhost:3000/produtos} $\rightarrow$ deve retornar a lista.
  \item \textbf{POST} \texttt{http://localhost:3000/produtos} com corpo JSON:
\begin{lstlisting}[language=JavaScript]
{ "nome": "Régua 30cm", "preco": 4.50, "estoque": 30 }
\end{lstlisting}
  \item \textbf{PUT} \texttt{http://localhost:3000/produtos/1} com novos dados.
  \item \textbf{DELETE} \texttt{http://localhost:3000/produtos/1}.
\end{enumerate}

\begin{caixa-pratica}
Crie 3 produtos novos via POST e confirme que aparecem no GET. Atualize um deles e depois apague.
\end{caixa-pratica}

% ============================================================
\section{Passo 8 --- HTML do frontend}
% ============================================================

\begin{lstlisting}[language=HTML, caption={\texttt{frontend/index.html}}]
<!DOCTYPE html>
<html lang="pt-BR">
<head>
  <meta charset="UTF-8">
  <title>CRUD de Produtos</title>
  <link rel="stylesheet" href="style.css">
</head>
<body>
  <h1>Cadastro de Produtos</h1>

  <form id="form-produto">
    <input type="hidden" id="produto-id">
    <label>Nome: <input type="text" id="nome" required></label>
    <label>Preco: <input type="number" id="preco" step="0.01" required></label>
    <label>Estoque: <input type="number" id="estoque" value="0" required></label>
    <div class="botoes">
      <button type="submit" id="btn-salvar">Salvar</button>
      <button type="button" id="btn-cancelar" hidden>Cancelar</button>
    </div>
  </form>

  <h2>Lista de Produtos</h2>
  <table id="tabela-produtos">
    <thead>
      <tr>
        <th>ID</th><th>Nome</th><th>Preco</th>
        <th>Estoque</th><th>Acoes</th>
      </tr>
    </thead>
    <tbody></tbody>
  </table>

  <script src="script.js"></script>
</body>
</html>
\end{lstlisting}

% ============================================================
\section{Passo 9 --- JavaScript do frontend}
% ============================================================

Usamos a \texttt{Fetch API} para conversar com o backend.

\begin{lstlisting}[language=JavaScript, caption={\texttt{frontend/script.js (parte 1)}}]
const API = 'http://localhost:3000/produtos';

const form = document.getElementById('form-produto');
const inputId = document.getElementById('produto-id');
const inputNome = document.getElementById('nome');
const inputPreco = document.getElementById('preco');
const inputEstoque = document.getElementById('estoque');
const btnCancelar = document.getElementById('btn-cancelar');
const tbody = document.querySelector('#tabela-produtos tbody');

function criarCelula(texto) {
  const td = document.createElement('td');
  td.textContent = texto;
  return td;
}

function criarBotao(rotulo, classe, id) {
  const botao = document.createElement('button');
  botao.textContent = rotulo;
  botao.className = classe;
  botao.dataset.id = id;
  return botao;
}
\end{lstlisting}

\begin{lstlisting}[language=JavaScript, caption={\texttt{script.js (parte 2 --- CRUD)}}]
async function listarProdutos() {
  const resposta = await fetch(API);
  const produtos = await resposta.json();
  tbody.replaceChildren();
  produtos.forEach((p) => {
    const tr = document.createElement('tr');
    tr.append(
      criarCelula(p.id),
      criarCelula(p.nome),
      criarCelula(`R$ ${Number(p.preco).toFixed(2)}`),
      criarCelula(p.estoque),
    );
    const tdAcoes = document.createElement('td');
    tdAcoes.append(
      criarBotao('Editar', 'btn-editar', p.id),
      criarBotao('Excluir', 'btn-excluir', p.id),
    );
    tr.append(tdAcoes);
    tbody.append(tr);
  });
}

async function salvarProduto(evento) {
  evento.preventDefault();
  const dados = {
    nome: inputNome.value,
    preco: parseFloat(inputPreco.value),
    estoque: parseInt(inputEstoque.value, 10),
  };
  const id = inputId.value;
  const url = id ? `${API}/${id}` : API;
  const metodo = id ? 'PUT' : 'POST';
  await fetch(url, {
    method: metodo,
    headers: { 'Content-Type': 'application/json' },
    body: JSON.stringify(dados),
  });
  resetarFormulario();
  listarProdutos();
}

async function editarProduto(id) {
  const resposta = await fetch(`${API}/${id}`);
  const produto = await resposta.json();
  inputId.value = produto.id;
  inputNome.value = produto.nome;
  inputPreco.value = produto.preco;
  inputEstoque.value = produto.estoque;
  btnCancelar.hidden = false;
}

async function excluirProduto(id) {
  if (!confirm('Deseja excluir este produto?')) return;
  await fetch(`${API}/${id}`, { method: 'DELETE' });
  listarProdutos();
}

function resetarFormulario() {
  form.reset();
  inputId.value = '';
  inputEstoque.value = '0';
  btnCancelar.hidden = true;
}

form.addEventListener('submit', salvarProduto);
btnCancelar.addEventListener('click', resetarFormulario);

tbody.addEventListener('click', (evento) => {
  const id = evento.target.dataset.id;
  if (!id) return;
  if (evento.target.classList.contains('btn-editar')) editarProduto(id);
  else if (evento.target.classList.contains('btn-excluir')) excluirProduto(id);
});

listarProdutos();
\end{lstlisting}

\begin{caixa-atencao}
Use \texttt{textContent} e \texttt{createElement} ao invés de \texttt{innerHTML} para evitar \textbf{XSS} (cross-site scripting) quando exibir dados vindos da API.
\end{caixa-atencao}

% ============================================================
\section{Passo 10 --- CSS mínimo}
% ============================================================

\begin{lstlisting}[language=CSS, caption={\texttt{frontend/style.css (resumido)}}]
* { box-sizing: border-box; }
body {
  font-family: system-ui, sans-serif;
  max-width: 800px; margin: 2rem auto; padding: 0 1rem;
  background: #f5f5f5; color: #222;
}
form, table {
  background: #fff; border-radius: 8px;
  box-shadow: 0 1px 3px rgba(0,0,0,.1);
}
form { padding: 1.5rem; margin-bottom: 2rem; }
form input { width: 100%; padding: .5rem; margin-top: .25rem; }
button { padding: .5rem 1rem; border: none; border-radius: 4px; cursor: pointer; }
#btn-salvar  { background: #2563eb; color: #fff; }
.btn-editar  { background: #f59e0b; color: #fff; }
.btn-excluir { background: #dc2626; color: #fff; }
table { width: 100%; border-collapse: collapse; }
th, td { padding: .75rem; text-align: left; border-bottom: 1px solid #eee; }
\end{lstlisting}

% ============================================================
\section{Passo 11 --- Rodar tudo junto}
% ============================================================

\begin{enumerate}[itemsep=2pt]
  \item Em um terminal: \texttt{cd backend \&\& bun run dev}.
  \item Abra \texttt{frontend/index.html} no navegador (botão direito $\rightarrow$ \emph{Open with Live Server} no VS Code é o ideal).
  \item Cadastre, edite e exclua produtos. Veja a lista atualizar.
\end{enumerate}

\begin{caixa-pratica}
Confira no DevTools (F12) $\rightarrow$ \emph{Network}: cada ação dispara uma requisição diferente para a API.
\end{caixa-pratica}

% ============================================================
\section{Resolução de problemas comuns}
% ============================================================

\begin{description}[leftmargin=*, itemsep=4pt]
  \item[\textcolor{vermelhoErr}{ERR\_CONNECTION\_REFUSED}] Backend não está rodando. Volte ao terminal e suba com \texttt{bun run dev}.
  \item[\textcolor{vermelhoErr}{CORS blocked}] Faltou \texttt{app.use(cors())} antes das rotas no \texttt{server.js}.
  \item[\textcolor{vermelhoErr}{ER\_ACCESS\_DENIED}] Usuário/senha errados em \texttt{config/database.js}.
  \item[\textcolor{vermelhoErr}{ER\_BAD\_DB\_ERROR}] Banco \texttt{dsw\_crud} não existe. Rode o \texttt{database.sql}.
  \item[\textcolor{vermelhoErr}{Cannot read properties of undefined}] Esqueceu \texttt{express.json()} --- \texttt{req.body} fica vazio.
  \item[\textcolor{vermelhoErr}{Tabela não atualiza}] Verifique se a função \texttt{listarProdutos()} é chamada após cada POST/PUT/DELETE.
\end{description}

% ============================================================
\section{Desafio final}
% ============================================================

\begin{caixa-pratica}
\textbf{Adicione um campo \texttt{categoria} (VARCHAR 50)} ao produto. Para isso você precisa:
\begin{enumerate}[itemsep=2pt]
  \item Alterar a tabela: \texttt{ALTER TABLE produtos ADD categoria VARCHAR(50);}
  \item Atualizar as queries de \texttt{INSERT} e \texttt{UPDATE} no controller.
  \item Adicionar o campo no formulário HTML.
  \item Adicionar o campo no objeto \texttt{dados} do \texttt{salvarProduto}.
  \item Adicionar a coluna na tabela do frontend.
\end{enumerate}
Repare como uma única feature toca \textbf{quatro camadas}: banco, controller, HTML e JS.
\end{caixa-pratica}

% ============================================================
\section{Checklist de entrega}
% ============================================================

\begin{itemize}[label=$\square$, itemsep=2pt]
  \item Banco \texttt{dsw\_crud} criado com a tabela \texttt{produtos}.
  \item Backend rodando em \texttt{localhost:3000} sem erros.
  \item Endpoint \texttt{GET /produtos} retorna JSON.
  \item Frontend lista, cria, edita e exclui produtos.
  \item Estrutura \texttt{routes/} + \texttt{controllers/} + \texttt{config/} respeitada.
  \item Campo \texttt{categoria} (desafio) implementado.
\end{itemize}

\vspace{0.6cm}
\hrule
\vspace{0.3cm}
\begin{center}\small
\textit{Bom código!} --- Prof. Ivan Pires \\
\texttt{ivan.pires@unemat.br} --- DSW 2026.1
\end{center}

\end{document}
